% ==============================================================================
% Section 1: Introduction
% ==============================================================================
\section{Introduction}
\label{sec:intro}

Public procurement accounts for approximately 12\% of GDP in OECD countries and an even larger share in developing economies, where government purchasing represents a critical channel for public service delivery \citep{hovenkamp2005federal}. Bid rigging---coordinated bidding among ostensibly competing firms---is among the most pervasive and costly forms of collusion in this setting. The OECD estimates that cartels inflate procurement prices by 10--50\%, with a median overcharge of approximately 20\% \citep{connor2007price}. Despite these substantial welfare losses, bid rigging remains notoriously difficult to detect. Competition authorities worldwide rely on a combination of leniency programs, whistleblowers, and---increasingly---quantitative screening tools to identify suspicious bidding patterns \citep{harrington2008detecting, imhof2019detecting}. However, most existing screens operate at the \textit{bid level}, flagging individual tenders based on statistical anomalies in bid distributions. These methods require substantial data on bid values, demand specialized econometric expertise, and produce high false-positive rates when applied to heterogeneous procurement environments where bid distributions naturally vary across product categories and time periods \citep{sanchezgraells2019screening}.

We propose a complementary approach: a \textit{firm-level} screen based on \textit{frequent losers} (FL)---firms that systematically participate in procurement tenders without ever winning. We argue that such firms are consistent with the role of \textit{cover bidders} in cartel arrangements: entities that submit non-competitive bids to satisfy minimum bidder requirements or create an illusion of competition, enabling the designated winner to secure the contract at an inflated price \citep{porter1993detection, marshall2012economics}. The operational logic is straightforward: in many procurement systems, a minimum number of bidders is required for the tender to proceed. A cartel that controls the designated winner must still populate the tender with enough bids to meet this threshold. Cover bidders---firms that participate with no intention of winning---serve this function. A firm that participates in dozens or hundreds of tenders without ever winning a single contract exhibits precisely the behavioral pattern predicted by the cover-bidding hypothesis.

The FL screen offers two practical advantages over existing bid-level screens. First, it requires only firm-level participation and outcome data (who bid, who won), not detailed bid values. This makes it deployable in procurement systems where bid prices are incompletely recorded or where data access is restricted. Second, because the FL classification operates at the firm level, a single flag identifies all tenders in which the suspect firm participates---potentially hundreds of tenders from a single classification. This yields a much higher screening throughput per unit of analytical effort compared with tender-by-tender statistical tests.

\subsection{Contributions}

This paper makes three contributions to the literature on cartel detection in procurement:

\begin{enumerate}
\item \textbf{Firm-level screening marker with external validation.} We propose frequent losers as a novel firm-level screen for cover bidding. Unlike bid-level screens \citep{bajari2003deciding, imhof2019detecting}, the FL marker operates at the firm level: a single classification flags all tenders in which the firm participates. This makes the screen operationally simpler for competition authorities. We validate the FL screen against CADE cartel convictions, showing that 193 FL firms co-participate in tenders with convicted cartelists, and that FL-present tenders exhibit 6--9\% higher winning prices across OLS, CEM matching, and IPW estimators. Standard Imhof (2019) bid-rigging screens achieve an AUC of 0.77 in predicting FL presence, indicating that the FL screen captures bid-level anomalies without requiring bid price data.

\item \textbf{Regime identification.} We empirically distinguish between two cover-bidding regimes: complementary cover bidding (Regime 1), where FL firms submit bids drawn independently above the reference price; and coordinated cover bidding (Regime 2), where FL bids cluster near the winning price to simulate competition. We find that FL bid dispersion is significantly higher than non-FL bid dispersion, consistent with Regime~1. Additionally, 91\% of FL bids exceed the winning price (median spread: 55\%), suggesting that FL firms submit transparently non-competitive bids rather than coordinating to mimic legitimate competition.

\item \textbf{Statistical evidence of coordination.} \citet{bajari2003deciding} tests reject both exchangeability and conditional independence for FL bid residuals: the distribution of FL bids differs significantly from non-FL bids (KS statistic = 0.15, $p < 0.001$), and FL bid residuals within the same tender are positively correlated (mean pairwise product = 5.16, $t = 157.4$). The baseline correlation among non-FL bids (mean product = 2.21) is smaller, though also statistically significant---consistent with common cost shocks affecting all bidders, with an \textit{additional} coordination component among FL firms.
\end{enumerate}

We are transparent about the limits of our identification strategy. The FL-price association is robust to matching estimators, sensitivity analysis, and threshold variation, but we cannot definitively establish that FL presence \textit{causes} higher prices rather than being correlated with unobserved market characteristics that independently drive prices upward. Staggered difference-in-differences analysis, which could in principle address this concern, does not yield clean pre-trends in our setting. We therefore frame the FL screen as a \textit{screening marker}---a statistical pattern strongly consistent with cover bidding and externally validated against known cartels---rather than as definitive proof of cartel activity in every flagged tender.

The remainder of this paper is organized as follows. Section~\ref{sec:literature} reviews the related literature. Section~\ref{sec:theory} presents the theoretical framework of cover bidding equilibrium with testable predictions. Section~\ref{sec:data} describes the data and the frequent loser definition. Section~\ref{sec:cade} presents external validation against CADE cartel convictions. Section~\ref{sec:results} reports the main results, including Bajari-Ye tests and the regime test. Section~\ref{sec:mechanisms} discusses mechanisms and alternative explanations. Section~\ref{sec:robustness} summarizes robustness checks and extensions, with full details in the Appendix. Section~\ref{sec:conclusion} concludes.
